\documentclass[12pt,a4paper]{article}

% ----- PACKAGES -----
\usepackage[utf8]{inputenc}
\usepackage[T1]{fontenc}
\usepackage[french]{babel}
\usepackage{geometry}
\geometry{margin=2.5cm}
\usepackage{graphicx}
\usepackage{hyperref}
\usepackage{listings}
\usepackage{xcolor}
\usepackage{amsmath}
\usepackage{amssymb}
\usepackage{array}
\usepackage{tabularx}
\usepackage{booktabs}
\usepackage{caption}
\usepackage{enumitem}
\usepackage{fancyhdr}
\usepackage{titlesec}

% ----- CODE LISTINGS STYLE -----
\definecolor{codebg}{rgb}{0.95,0.95,0.95}
\definecolor{codegreen}{rgb}{0,0.5,0}
\definecolor{codeblue}{rgb}{0,0,0.7}
\definecolor{codered}{rgb}{0.7,0,0}

\lstdefinestyle{mystyle}{
    backgroundcolor=\color{codebg},
    commentstyle=\color{codegreen},
    keywordstyle=\color{codeblue},
    stringstyle=\color{codered},
    basicstyle=\ttfamily\small,
    breakatwhitespace=false,
    breaklines=true,
    captionpos=b,
    keepspaces=true,
    showspaces=false,
    showstringspaces=false,
    showtabs=false,
    tabsize=2,
    frame=single,
    numbers=left,
    numberstyle=\tiny\color{gray}
}
\lstset{style=mystyle}

% ----- HEADER / FOOTER -----
\pagestyle{fancy}
\fancyhf{}
\fancyhead[L]{Pipeline IoT pour Raffinerie}
\fancyhead[R]{\thepage}
\fancyfoot[C]{}

% ----- TITLE -----
\title{\textbf{Pipeline IoT Industrielle pour la Surveillance Prédictive d'une Raffinerie}\\[1em]
\large Architecture Big Data --- Génération de Données --- Inférence ML Temps Réel}
\author{\textit{Document de préparation à l'entretien technique}}
\date{\today}

% ----- DOCUMENT -----
\begin{document}

\maketitle
\thispagestyle{empty}
\newpage

\tableofcontents
\newpage

% ============================================================
% PARTIE 1 : PIPELINE COMPLETE
% ============================================================
\section{Architecture Globale de la Pipeline}
\label{sec:pipeline}

\subsection{Vue d'ensemble}

Le projet met en œuvre une pipeline Big Data de bout en bout pour la surveillance prédictive d'équipements industriels (échangeur thermique \texttt{pipe-101}) dans une raffinerie. Les données simulées transitent par une architecture en couches :

\begin{center}
\small
\begin{tabular}{|c|c|c|}
\hline
\textbf{Couche} & \textbf{Technologie} & \textbf{Rôle} \\
\hline
Acquisition & Simulateur Python + MQTT & Génération de données capteurs réalistes \\
\hline
Transport & Kafka & File d'attente distribué, tampon de messages \\
\hline
Traitement & Spark Structured Streaming & Filtrage, KPI, inférence ML en continu \\
\hline
Stockage long terme & TimescaleDB (hypertables) & Séries temporelles, prédictions \\
\hline
Stockage froid & MinIO (S3) & Archivage, données d'entraînement \\
\hline
Visualisation & Grafana & Dashboard temps réel, jauges, alertes \\
\hline
\end{tabular}
\end{center}

\subsection{Flux des données pas à pas}

Le cheminement d'une mesure, de sa création à son affichage, est le suivant :

\begin{enumerate}
    \item \textbf{Simulateur} (\texttt{simulateur\_capteurs.py}) : génère une paire température + vibration toutes les 2~secondes avec une logique physique (inertie, bruit, scénarios aléatoires).

    \item \textbf{MQTT} (\texttt{eclipse-mosquitto}) : le simulateur publie sur deux topics --- \texttt{raffinerie/temp} et \texttt{raffinerie/vib} --- via le protocole MQTT (léger, adapté à l'IoT).

    \item \textbf{Pont MQTT→Kafka} (\texttt{mqtt\_to\_kafka.py}) : un bridge s'abonne aux topics MQTT et publie chaque message dans un topic Kafka unique \texttt{sensor-data}. Ce découplage permet à Kafka d'absorber les pics de charge.

    \item \textbf{Kafka} (\texttt{confluentinc/cp-kafka}) : joue le rôle de tampon distribué. Les messages sont persistés et réplicables, garantissant qu'aucune donnée n'est perdue même si Spark est momentanément indisponible.

    \item \textbf{Spark Structured Streaming} (\texttt{traitement\_kpi.py}) : lit le topic Kafka en continu et exécute trois traitements en parallèle :
    \begin{itemize}
        \item Filtrage des valeurs aberrantes (température 30--150\,°C, vibration 0--5\,mm/s) et écriture dans \texttt{mesures\_filtrees}.
        \item Calcul de KPI (moyennes glissantes 1~minute) écrits dans \texttt{kpi\_indicateurs}.
        \item \textbf{Inférence ML} : application du modèle RandomForest pour produire un score d'anomalie, écrit dans \texttt{alertes\_predictions}.
    \end{itemize}

    \item \textbf{TimescaleDB} (\texttt{timescale/timescaledb:2.14.2-pg14}) : base de données spécialisée séries temporelles (hypertables). Les index DESC sur le timestamp optimisent les requêtes \texttt{ORDER BY timestamp DESC LIMIT N}.

    \item \textbf{Grafana} (\texttt{grafana/grafana:10.2.3}) : dashboards rafraîchis toutes les 5~secondes interrogeant directement TimescaleDB via un datasource PostgreSQL.
\end{enumerate}

\subsection{Déploiement Docker}

Les 11 services sont déployés via Docker Compose sur un réseau bridge \texttt{iot-net} :

\begin{lstlisting}[language=yaml,caption={Structure du docker-compose.yml},label={lst:compose}]
services:
  mqtt:         image: eclipse-mosquitto:2.0        ports: 1883
  zookeeper:    image: confluentinc/cp-zookeeper:7.4.0
  kafka:        image: confluentinc/cp-kafka:7.4.0  ports: 9092, 29092
  spark-master: image: apache/spark:3.4.1           ports: 7077, 8080
  spark-worker: image: apache/spark:3.4.1
  minio:        image: minio/minio                   ports: 9000, 9001
  timescaledb:  image: timescale/timescaledb:2.14.2-pg14  ports: 5432
  grafana:      image: grafana/grafana:10.2.3        ports: 3000
  simulateur:   build: ./simulateur                  # Multi-machines, lecture DB
  mqtt-bridge:  build: ./bridge                      # Wildcard topics MQTT
  webapp:       build: ./webapp                      # FastAPI + HTMX, ports: 8081
\end{lstlisting}

\newpage

\subsection{Détail des volumes et persistances}

\begin{itemize}
    \item \textbf{\texttt{./spark:/app}} (master uniquement) : contient les scripts Python (\texttt{traitement\_kpi.py}, \texttt{train\_model.py}).
    \item \textbf{\texttt{spark\_ivy:/home/spark/.ivy2}} (master + worker) : volume partagé contenant le cache Ivy (dépendances Spark), numpy, ainsi que le \textbf{modèle prédictif} copié dans \texttt{/home/spark/.ivy2/predictive\_model}.
    \item \textbf{\texttt{minio\_data:/data}} : stockage objet MinIO (données brutes, données d'entraînement).
    \item \textbf{\texttt{pgdata:/var/lib/postgresql/data}} : données TimescaleDB.
    \item \textbf{\texttt{grafana\_data:/var/lib/grafana}} : configuration et métadonnées Grafana.
\end{itemize}

\subsection{Emplacement des données et artefacts critiques}

\begin{center}
\small
\begin{tabularx}{\textwidth}{|l|X|c|}
\hline
\textbf{Quoi} & \textbf{Où} & \textbf{Taille} \\
\hline
Données d'entraînement labellisées & MinIO : \texttt{raffinerie-raw/training\_data/training\_data.json} & 10\,000 enreg. \\
\hline
Modèle ML (PipelineModel 3 stages) & \texttt{/home/spark/.ivy2/predictive\_model/} & 3 stages, 30 arbres \\
\hline
Mesures filtrées temps réel & TimescaleDB : \texttt{mesures\_filtrees} (hypertable) & $\sim$200\,000 lignes \\
\hline
KPI (moyennes 1~min) & TimescaleDB : \texttt{kpi\_indicateurs} (hypertable) & $\sim$6\,000 lignes \\
\hline
Prédictions / alertes & TimescaleDB : \texttt{alertes\_predictions} (hypertable) & $\sim$72\,000 lignes \\
\hline
Checkpoints Spark & \texttt{/tmp/checkpoint\_*} (dans le conteneur master) & Volatil \\
\hline
\end{tabularx}
\end{center}

\subsection{Orchestration et lancement}

Le lancement est entièrement géré par Docker Compose :

\begin{itemize}
    \item \texttt{docker compose up -d} : démarre les 11 services (MQTT, Kafka, Zookeeper, Spark master/worker, MinIO, TimescaleDB, Grafana, simulateur, bridge MQTT-Kafka, webapp). L'ordre de démarrage est contrôlé par \texttt{depends\_on}.
    \item \texttt{restart.sh} : script de commodité qui exécute \texttt{docker compose up -d --build} et vérifie que tous les conteneurs sont opérationnels.
    \item \textbf{Webapp} : accessible sur \texttt{http://localhost:8081}, permet de lancer/arrêter la pipeline (simulateur + bridge + Spark) via des boutons HTMX. Un lien direct vers Grafana est présent dans la barre de navigation.
    \item \textbf{Gestion des machines} : depuis l'interface web, l'opérateur peut ajouter, supprimer, activer ou désactiver des machines. Le simulateur détecte automatiquement les changements dans les 60~secondes.
    \item \texttt{spark-submit} : soumet \texttt{traitement\_kpi.py} sur le cluster Spark (master:7077). Les paquets Kafka et PostgreSQL sont injectés via \texttt{--packages}. Le lancement se fait automatiquement via l'interface web ou manuellement par \texttt{docker exec spark-master spark-submit ...}.
    \item \texttt{spark-submit} pour l'entraînement : \texttt{train\_model.py} lit depuis MinIO (via S3A Hadoop), entraîne, et sauvegarde localement.
\end{itemize}

\subsection{Configuration Spark particulière}

Le worker Spark n'a pas accès au volume \texttt{./spark:/app}. Deux ajustements ont été nécessaires :

\begin{enumerate}
    \item \textbf{Modèle sur volume partagé} : \texttt{PipelineModel.load("/home/spark/.ivy2/predictive\_model")} (le volume \texttt{spark\_ivy} est monté sur les deux nœuds).
    \item \textbf{NumPy pour executors} : installé dans \texttt{/home/spark/.ivy2/packages}, accessible via \texttt{spark.executorEnv.PYTHONPATH}.
\end{enumerate}

\newpage

% ============================================================
% PARTIE 2 : AXES TECHNIQUES
% ============================================================
\section{Axes de réalisation}

\subsection{Axe 1 : Génération de données réalistes par les capteurs IoT}

\subsubsection{Objectif}

Produire un flux de données synthétiques imitant le comportement physique de capteurs industriels (température, vibration) avec des dérives réalistes, afin d'entraîner et valider un modèle de détection d'anomalies.

\subsubsection{Logique physique : lissage exponentiel avec inertie}

Un capteur réel ne change pas instantanément. La valeur $V_t$ à l'instant $t$ dépend de la valeur précédente $V_{t-1}$, d'une valeur cible $Cible$, d'un facteur d'inertie $\alpha$, et d'un bruit $\varepsilon$ :

\[
V_t = V_{t-1} \times (1 - \alpha) + Cible \times \alpha + \varepsilon
\]

\begin{itemize}
    \item $\alpha_{\text{temp}} = 0.05$ : inertie forte (la température met $\sim$20~itérations à converger).
    \item $\alpha_{\text{vib}} = 0.2$ : inertie modérée (les vibrations réagissent plus vite).
    \item $\varepsilon_{\text{temp}} \sim \mathcal{U}(-0.2, 0.2)$ ; $\varepsilon_{\text{vib}} \sim \mathcal{U}(-0.05, 0.05)$.
\end{itemize}

Cette formule modélise à la fois l'inertie thermique d'un échangeur et l'inertie mécanique d'une pompe.

\subsubsection{Scénarios stochastiques}

Toutes les 30 à 60 itérations, un dé à trois faces décide du prochain scénario :

\begin{center}
\begin{tabular}{|l|c|c|c|}
\hline
\textbf{Scénario} & \textbf{Probabilité} & \textbf{Cible Temp.} & \textbf{Cible Vib.} \\
\hline
NORMAL & 90\% & 95.0\,°C & 1.2\,mm/s \\
SURCHAUFFE & 5\% & 145.0\,°C & 1.2\,mm/s \\
USURE MÉCANIQUE & 5\% & 95.0\,°C & 4.5\,mm/s \\
\hline
\end{tabular}
\end{center}

En mode SURCHAUFFE, seule la température dérive (la pompe reste normale). En mode USURE, seules les vibrations augmentent. Ce cloisonnement force le modèle à apprendre les deux types d'anomalie indépendamment.

\subsubsection{Contraintes de clipping}

Pour respecter les filtres aval de Spark (30--150\,°C, 0--5\,mm/s), les valeurs sont clampées :

\[
\text{temp} = \max(31, \min(149, \text{temp}))
\]
\[
\text{vib} = \max(0.1, \min(4.9, \text{vib}))
\]

\subsubsection{Publication MQTT}

Chaque itération produit deux messages JSON (température + vibration) avec le même \texttt{timestamp}, publiés sur les topics MQTT \texttt{raffinerie/temp} et \texttt{raffinerie/vib}.

\begin{lstlisting}[language=Python,caption={Structure d'un message MQTT},label={lst:mqttmsg}]
{
  "machine_id": "pipe-101",
  "valeur": 95.08,
  "timestamp": "2026-05-18T20:40:00Z",
  "type_capteur": "temperature"
}
\end{lstlisting}

\subsubsection{Scripts associés}

\begin{itemize}
    \item \texttt{simulateur\_capteurs.py} : le simulateur lui-même, tourne en boucle infinie avec \texttt{time.sleep(2)}.
    \item \texttt{mqtt\_to\_kafka.py} : pont MQTT→Kafka, s'abonne aux deux topics et forwarde dans Kafka topic \texttt{sensor-data}.
    \item \texttt{LOGIQUE\_SIMULATION.md} : documentation détaillée des principes physiques (inertie, corrélation, scénarios).
\end{itemize}

\newpage

\subsection{Axe 2 : Prédiction intelligente et détection d'anomalies}

\subsubsection{Génération des données d'entraînement}

Le script \texttt{generate\_training\_data.py} reproduit exactement la même logique physique que le simulateur, mais en mode batch (5\,000 itérations $\rightarrow$ 10\,000 enregistrements). Chaque enregistrement est labellisé :

\[
\text{label} = \begin{cases}
0 & \text{si scénario = NORMAL} \\
1 & \text{si scénario = SURCHAUFFE ou USURE\_MECANIQUE}
\end{cases}
\]

Les données sont uploadées sur MinIO au format JSON Lines :

\[
\text{MinIO} : \texttt{raffinerie-raw/training\_data/training\_data.json}
\]

\subsubsection{Données d'entraînement : distribution et split}

Le dataset généré contient 10\,000 enregistrements (5\,000 itérations $\times$ 2 capteurs) avec la répartition suivante :

\begin{center}
\begin{tabular}{|c|c|c|}
\hline
\textbf{Label} & \textbf{Effectif} & \textbf{Pourcentage} \\
\hline
0 (NORMAL) & 9\,140 & 91,4\% \\
1 (ANOMALIE) & 860 & 8,6\% \\
\hline
\end{tabular}
\end{center}

Cette distribution reflète les probabilités des scénarios (90\% normal, 5\% surchauffe, 5\% usure). Les données sont stockées sur MinIO au format JSON Lines :

\[
\texttt{s3a://raffinerie-raw/training\_data/training\_data.json}
\]

Le jeu est ensuite divisé aléatoirement (80\% train, 20\% test) via \texttt{randomSplit([0.8, 0.2], seed=42)}.

\subsubsection{Feature engineering}

Chaque mesure brute $(valeur, type\_capteur, timestamp)$ est transformée en un vecteur de 9 features :

\begin{enumerate}
    \item \textbf{Valeur instantanée} : la mesure elle-même ($V_t$).
    \item \textbf{Lags temporels} ($\times 5$) : $V_{t-1}, V_{t-2}, V_{t-3}, V_{t-4}, V_{t-5}$ — soit une fenêtre de 10~secondes (5 pas $\times$ 2~s).
    \item \textbf{Pentes} ($\times 2$) : $\text{slope\_1} = V_t - V_{t-1}$ (pente immédiate), $\text{slope\_3} = (V_t - V_{t-3}) / 3$ (tendance sur 6~s). Ces pentes sont les \textbf{features les plus importantes} pour la détection précoce : une valeur encore normale mais en hausse produit un score d'anomalie élevé.
    \item \textbf{Type de capteur} ($\times 1$) : encodé via \texttt{StringIndexer} (temperature $\rightarrow$ 0.0, vibration $\rightarrow$ 1.0).
\end{enumerate}

Les lags sont calculés par fenêtre de partition Spark :

\[
\texttt{Window.partitionBy("machine\_id", "type\_capteur").orderBy("timestamp")}
\]

Les lignes avec des lags nuls (début de séquence, pas d'historique) sont supprimées via \texttt{dropna()}.

\subsubsection{Sélection de l'algorithme : pourquoi RandomForest ?}

Le choix du RandomForest (plutôt qu'un XGBoost, un LSTM ou un SVM) repose sur plusieurs critères adaptés à ce contexte industriel :

\begin{itemize}
    \item \textbf{Passage à l'échelle natif dans Spark} : \texttt{RandomForestClassifier} est implémenté directement dans MLlib, sans dépendance externe. L'entraînement et l'inférence se font \textbf{dans le même framework} que le streaming, évitant un déploiement séparé.
    \item \textbf{Robustesse au déséquilibre de classe} (91\% / 9\%) : le vote majoritaire des arbres et le bagging intrinsèque gèrent le déséquilibre sans nécessiter de sur-échantillonnage ou de pondération explicite.
    \item \textbf{Interprétabilité} : l'importance des features est accessible, et on peut expliquer pourquoi une prédiction est positive (pente anormale). Pas de boîte noire.
    \item \textbf{Performance en temps réel} : l'inférence d'une forêt de 30 arbres depth 6 sur 9 features prend quelques millisecondes par ligne, compatible avec le trigger à 2~secondes.
    \item \textbf{Données tabulaires structurées} : le problème est un classifieur binaire sur features numériques et catégorielle, pas une séquence longue. Un LSTM serait surdimensionné et coûteux pour un gain marginal.
\end{itemize}

\subsubsection{Pipeline d'entraînement}

Le pipeline Spark ML (3 stages) est assemblé et exécuté en une ligne :

\begin{lstlisting}[language=Python,caption={Pipeline d'entraînement (train\_model.py)},label={lst:pipeline}]
indexer = StringIndexer(inputCol="type_capteur", outputCol="sensor_idx")
assembler = VectorAssembler(
    inputCols=["valeur", "v_lag1", ..., "v_lag5",
               "slope_1", "slope_3", "sensor_idx"],
    outputCol="features")
rf = RandomForestClassifier(labelCol="label", featuresCol="features",
                            numTrees=30, maxDepth=6, seed=42)
pipeline = Pipeline(stages=[indexer, assembler, rf])
model = pipeline.fit(train_data)
\end{lstlisting}

\subsubsection{Métriques et évaluation}

Le modèle est évalué sur le jeu de test (20\%, soit $\sim$2\,000 lignes) avec un \texttt{BinaryClassificationEvaluator} :

\[
\text{AUC} = 0.8027
\]

L'AUC (Area Under the ROC Curve) mesure la capacité du modèle à séparer les deux classes, indépendamment du seuil. Un AUC de 0.80 signifie qu'il y a 80\% de chances que le modèle classe correctement une paire (normal, anomalie) aléatoire. C'est un score acceptable pour un premier entraînement sans optimisation d'hyperparamètres (grid search). Les axes d'amélioration possibles seraient : augmenter \texttt{numTrees} (100+), réduire \texttt{maxDepth} (5) pour limiter le sur-apprentissage, ou appliquer du SMOTE pour renforcer la classe minoritaire.

Le modèle final (30 arbres, depth 6, seed 42) est sauvegardé dans :

\[
\texttt{/home/spark/.ivy2/predictive\_model}
\]

Structure du modèle (3 stages) :
\begin{lstlisting}[language=bash]
predictive_model/
  metadata/                # Metadonnees du PipelineModel
  stages/
    0_StringIndexer_*/     # Mapping type_capteur -> sensor_idx
    1_VectorAssembler_*/   # Assemblage des 9 features
    2_RandomForestClassifier_*/  # 30 arbres, 2 classes
\end{lstlisting}

\newpage

\subsubsection{Inférence temps réel dans Spark Streaming}

La fonction \texttt{predict\_and\_alert} est exécutée via \texttt{foreachBatch} toutes les 2~secondes (trigger configurable) :

\begin{enumerate}
    \item Construction des mêmes features (lags + pentes) que lors de l'entraînement.
    \item \textbf{Gestion des lags nuls} : les premières lignes de chaque micro-batch ont des lags nuls (pas d'historique dans le batch). On utilise \texttt{coalesce(col("v\_lag\_k"), col("valeur"))} pour remplacer les nuls par la valeur courante (pente = 0), évitant de perdre la ligne.
    \item Application du \texttt{PipelineModel.transform()} qui produit les colonnes \texttt{prediction} et \texttt{probability}.
    \item Extraction de la probabilité de la classe 1 (anomalie) via une UDF :
    \[
    \text{score\_anomalie} = \text{probability}[1]
    \]
    \item Écriture dans \texttt{alertes\_predictions} (toutes les prédictions, pas seulement les dépassements de seuil).
\end{enumerate}

\subsubsection{Seuils et interprétation}

\begin{center}
\begin{tabular}{|c|c|c|}
\hline
\textbf{Score anomalie} & \textbf{Couleur Grafana} & \textbf{Interprétation} \\
\hline
$[0, 0.4[$ & Vert & Fonctionnement nominal \\
$[0.4, 0.7[$ & Jaune & Dérive détectée, surveillance renforcée \\
$[0.7, 1.0]$ & Rouge & Alerte critique, intervention requise \\
\hline
\end{tabular}
\end{center}

En fonctionnement normal, les scores sont :
\begin{itemize}
    \item Température stable à 95\,°C : $\sim$0.026
    \item Vibration stable à 1.2\,mm/s : $\sim$0.060
\end{itemize}

En anomalie, le score monte à $>$0.7 dès que la pente positive est détectée (2 à 4~secondes après le début de la dérive).

\subsubsection{Pourquoi les scores plafonnent-ils souvent à \textasciitilde0.7 ?}

Le RandomForestClassifier calcule la probabilité comme la proportion d'arbres ayant voté \og anomalie \fg{} (classe 1). Avec 30 arbres, les valeurs possibles sont des multiples de $1/30 \approx 0.033$.

Quand 21~arbres sur 30 détectent l'anomalie (0.7) mais que 9~votent \og normal \fg{} à cause de similarités résiduelles avec le régime normal (bruit, oscillations), le score se stabilise à 0.7. Atteindre 0.8 nécessiterait 24~arbres d'accord — possible mais plus rare.

\begin{center}
\begin{tabular}{|c|c|}
\hline
\textbf{Score} & \textbf{Signification} \\
\hline
0.02 -- 0.06 & Tous les arbres disent normal \\
0.3 -- 0.5 & Divergence naissante \\
\textbf{0.7} & \textbf{Majorité solide → alerte rouge} \\
0.9+ & Consensus quasi-unanime (rare à cause du bruit) \\
\hline
\end{tabular}
\end{center}

\noindent\textbf{L'important n'est pas le score absolu mais la séparation :} en fonctionnement normal, le score ne dépasse jamais 0.1. En anomalie, il dépasse toujours 0.7. Cette discontinuité garantit une détection fiable, quel que soit le plateau exact atteint (0.7 ou 0.8). Un calibrage probabiliste (Platt Scaling) serait possible mais n'apporterait rien ici, car la décision opérationnelle reste binaire : alerte déclenchée ou non.

\subsubsection{Comportement attendu par type d'anomalie}

\begin{itemize}
    \item \textbf{Surchauffe} (145\,°C) : la température monte avec une pente positive. Le modèle voit la pente de température (slope\_1 > 0) et score $>$0.7. La vibration reste normale ($\sim$0.06).
    \item \textbf{Usure mécanique} (4.5\,mm/s) : la vibration monte. Le modèle voit la pente de vibration et score $>$0.7. La température reste normale ($\sim$0.026).
\end{itemize}

\subsection{Visualisation Grafana}

\subsubsection{Panels du dashboard \og Surveillance Raffinerie \fg}

Le dashboard contient 4 panels :

\begin{enumerate}
    \item \textbf{Température par machine} : time series depuis \texttt{mesures\_filtrees} où \texttt{type\_capteur = 'temperature'}.
    \item \textbf{Vibrations par machine} : time series depuis \texttt{mesures\_filtrees} où \texttt{type\_capteur = 'vibration'}.
    \item \textbf{Alertes Prédictives -- Score d'Anomalie} : time series depuis \texttt{alertes\_predictions}, avec un seuil horizontal à 0.7 (ligne rouge).
    \item \textbf{Score d'Anomalie} : stat panel affichant la dernière valeur de \texttt{alertes\_predictions} avec un fond coloré (vert $<$ 0.4, jaune [0.4, 0.7[, rouge $\geq$ 0.7). La couleur donne un diagnostic immédiat sans interpréter le score brut.
\end{enumerate}

\subsubsection{Interprétation du score d'anomalie}

Le score affiché (\textit{stat panel} \og Niveau d'anomalie \fg) est la proportion d'arbres du RandomForest ayant voté \og anomalie \fg{} --- ce n'est \textbf{pas un pourcentage de risque}. En régime normal, 1 arbre sur 30 peut voter anomalie à cause du bruit, donnant $1/30 \approx 0.033$. La valeur $0.026$ observée en pratique signifie qu'en moyenne moins d'un arbre détecte une variation, ce qui est négligeable.

\begin{itemize}
    \item \textbf{0.000 -- 0.399} (fond \textcolor{green}{vert}) : fonctionnement nominal. Aucune action requise.
    \item \textbf{0.400 -- 0.699} (fond \textcolor{#d29922}{jaune}) : dérive détectée. Surveillance renforcée conseillée.
    \item \textbf{0.700 -- 1.000} (fond \textcolor{red}{rouge}) : anomalie confirmée. Intervention requise.
\end{itemize}

\textbf{Le score n'est pas une mesure de santé globale.} Il ne reflète que les deux capteurs simulés (température et vibration) et ne remplace pas une inspection réelle. Un score de 0.026 n'est pas un risque de 2.6\% --- c'est simplement la sortie du modèle quand tout est stable.

\subsubsection{Requêtes et rafraîchissement}

Chaque panel utilise une requête SQL directe vers le datasource PostgreSQL (TimescaleDB). Le dashboard est configuré en auto-refresh toutes les 5~secondes, mais les données arrivent toutes les 2~secondes via Spark, donc le facteur limitant est le rafraîchissement Grafana.

\newpage

% ============================================================
% AXE 3 : INTERFACE WEB ET DOCKERISATION
% ============================================================
\section{Axe 3 : Interface de Contrôle Web et Conteneurisation}
\label{sec:web}

\subsection{Objectif}

Permettre le contrôle et la configuration de la pipeline IoT depuis une interface web, sans nécessiter d'accès SSH ni de commandes manuelles. Les fonctionnalités attendues sont :
\begin{itemize}
    \item Lancer / arrêter / redémarrer la pipeline (simulateur, bridge, Spark).
    \item Configurer les machines simulées (ajout, suppression, paramètres cibles).
    \item Définir le seuil d'alerte du modèle de prédiction.
    \item Visualiser en temps réel l'état des services.
\end{itemize}

\subsection{Architecture de l'interface web}

L'interface utilise \textbf{FastAPI} (framework Python asynchrone) pour le backend et \textbf{HTMX} pour le frontend. Le choix HTMX plutôt qu'un framework JavaScript (React, Vue) repose sur plusieurs critères :

\begin{itemize}
    \item \textbf{Absence de build step} : pas de transpilation, pas de bundle. Les templates Jinja2 sont servis directement par le backend.
    \item \textbf{Architecture server-driven} : chaque action (start/stop, modification de config) provoque une requête HTTP vers une route FastAPI qui renvoie du HTML partiel.
    \item \textbf{3 écrans seulement} : dashboard, configuration, fragment de statut. Le surcoût d'un SPA serait disproportionné.
    \item \textbf{Mise à jour automatique} : HTMX permet de rafraîchir un fragment de page (\texttt{<div>}) toutes les 5~secondes via l'attribut \texttt{hx-trigger="every 5s"}.
\end{itemize}

\subsection{Structure de l'application}

\begin{lstlisting}[language=bash,caption={Arborescence de l'interface web}]
webapp/
  Dockerfile              # Image Python 3.11-slim
  requirements.txt        # fastapi, uvicorn, jinja2, docker, psycopg2-binary
  main.py                 # Routes FastAPI (7 routes)
  templates/
    base.html             # Layout commun (nav, style, HTMX CDN)
    dashboard.html        # Page d'accueil : statuts + controles
    config.html           # Configuration machines + seuil
    status_fragment.html  # Fragment HTMX (statut des services)
\end{lstlisting}

\subsection{Routes FastAPI}

\begin{center}
\small
\begin{tabular}{|l|l|l|}
\hline
\textbf{Méthode} & \textbf{Path} & \textbf{Fonction} \\
\hline
GET & \texttt{/} & Dashboard principal : statuts, dernière prédiction, nb machines \\
GET & \texttt{/config} & Page de configuration : liste machines + seuil \\
GET & \texttt{/status} & Fragment HTML des statuts (HTMX, refresh 5s) \\
POST & \texttt{/pipeline/start} & Démarre simulateur, bridge, soumet Spark \\
POST & \texttt{/pipeline/stop} & Arrête simulateur, bridge, tue Spark \\
POST & \texttt{/pipeline/restart} & Stop puis start (avec pause de 2s) \\
POST & \texttt{/config/machines/add} & Ajoute/met à jour une machine \\
POST & \texttt{/config/machines/delete/\{id\}} & Supprime une machine \\
POST & \texttt{/config/machines/toggle/\{id\}} & Active/désactive une machine \\
POST & \texttt{/config/threshold} & Modifie le seuil d'alerte \\
\hline
\end{tabular}
\end{center}

\subsection{Contrôle Docker depuis l'interface}

La webapp contrôle les conteneurs via le SDK Docker Python (\texttt{docker.from\_env()}) : le socket Docker (\texttt{/var/run/docker.sock}) est monté dans le conteneur webapp. Pour lancer le \texttt{spark-submit} dans le conteneur \texttt{spark-master}, la webapp utilise \texttt{container.exec\_run()} (détaché pour ne pas bloquer la réponse HTTP).

Les formulaires de configuration (ajout/suppression/toggle machine, seuil) renvoient un header \texttt{HX-Refresh: true} : HTMX rafraîchit la page automatiquement après l'action, évitant l'affichage de JSON brut.

\subsection{Configuration des machines et du seuil}

\subsubsection{Tables de configuration dans TimescaleDB}

Deux nouvelles tables stockent la configuration :

\begin{lstlisting}[language=SQL,caption={Tables de configuration}]
-- Machine config
CREATE TABLE machine_config (
    machine_id   TEXT PRIMARY KEY,
    machine_type TEXT NOT NULL DEFAULT 'pipe',
    enabled      BOOLEAN DEFAULT true,
    target_temp  DOUBLE PRECISION DEFAULT 95.0,
    target_vib   DOUBLE PRECISION DEFAULT 1.2,
    alpha_temp   DOUBLE PRECISION DEFAULT 0.05,
    alpha_vib    DOUBLE PRECISION DEFAULT 0.2,
    created_at   TIMESTAMPTZ DEFAULT now(),
    updated_at   TIMESTAMPTZ DEFAULT now()
);

-- Alert threshold config
CREATE TABLE alert_config (
    seuil_anomalie DOUBLE PRECISION DEFAULT 0.7,
    machine_id TEXT DEFAULT NULL,
    type_capteur TEXT DEFAULT NULL,
    updated_at TIMESTAMPTZ DEFAULT now()
);
\end{lstlisting}

\subsubsection{Propagation de la configuration}

La configuration est propagée dans la pipeline selon deux stratégies :

\begin{itemize}
    \item \textbf{Simulateur} : lit \texttt{machine\_config} au démarrage et toutes les 60~secondes. Si une machine est ajoutée ou activée, le simulateur commence à publier ses données. Si une machine est désactivée, elle est retirée du cycle de publication.
    \item \textbf{Spark} : lit \texttt{seuil\_anomalie} depuis \texttt{alert\_config} au démarrage du streaming. Le seuil n'est pas relu en continu (la modification du seuil nécessite un redémarrage de la pipeline via le bouton \og Redémarrer \fg{} de l'interface).
\end{itemize}

\subsection{Dockerisation des composants additionnels}

\subsubsection{Simulateur Dockerisé}

Le simulateur, auparavant script hôte lancé avec \texttt{nohup}, est maintenant un conteneur Docker :

\begin{lstlisting}[language=Dockerfile,caption={Dockerfile du simulateur}]
FROM python:3.11-slim
WORKDIR /app
RUN pip install paho-mqtt psycopg2-binary
COPY simulateur_capteurs.py .
ENV PYTHONUNBUFFERED=1
CMD ["python3", "simulateur_capteurs.py"]
\end{lstlisting}

Le simulateur Dockerisé lit sa configuration depuis TimescaleDB via \texttt{psycopg2}. Les topics MQTT sont désormais hiérarchiques :
\[
\texttt{raffinerie/\{machine\_id\}/temp} \quad \text{et} \quad \texttt{raffinerie/\{machine\_id\}/vib}
\]

Cette structure permet au bridge MQTT-Kafka de souscrire avec des wildcards :
\[
\texttt{raffinerie/+/temp} \quad \text{et} \quad \texttt{raffinerie/+/vib}
\]

Ainsi, ajouter une machine dans \texttt{machine\_config} suffit pour que le simulateur commence à publier sur son topic dédié, sans modifier aucun code ni configuration.

\subsubsection{Pont MQTT-Kafka Dockerisé}

Le bridge est lui aussi conteneurisé. Il intègre une boucle d'attente pour Kafka :

\begin{lstlisting}[language=Python,caption={Boucle d'attente pour Kafka dans mqtt\_to\_kafka.py}]
def wait_for_kafka(broker, retries=30, delay=2):
    for i in range(retries):
        try:
            p = KafkaProducer(bootstrap_servers=broker)
            p.close()
            return
        except NoBrokersAvailable:
            time.sleep(delay)
    sys.exit(1)  # Kafka inaccessible
\end{lstlisting}

Cette boucle résout le problème de \textit{race condition} au démarrage : Kafka peut prendre plusieurs secondes à être prêt, tandis que le bridge se lance immédiatement.

\subsection{Docker Compose mis à jour}

Le fichier \texttt{docker-compose.yml} passe de 8 à 11 services avec l'ajout de :

\begin{lstlisting}[language=yaml,caption={Nouveaux services dans docker-compose.yml}]
simulateur:
  build: ./simulateur
  depends_on: [mqtt, timescaledb]
  environment:
    MQTT_BROKER: mqtt
    DB_HOST: timescaledb

mqtt-bridge:
  build: ./bridge
  depends_on: [mqtt, kafka]
  environment:
    MQTT_BROKER: mqtt
    KAFKA_BROKER: kafka:29092

webapp:
  build: ./webapp
  ports: ["8081:8080"]
  volumes: ["/var/run/docker.sock:/var/run/docker.sock"]
  environment:
    DB_HOST: timescaledb
\end{lstlisting}

Le volume \texttt{initdb} est désormais monté dans TimescaleDB :
\[
\texttt{./initdb:/docker-entrypoint-initdb.d}
\]

Ce dossier contient \texttt{001\_config\_tables.sql} qui crée \texttt{machine\_config} et \texttt{alert\_config} avec les données par défaut au premier démarrage de la base.

\subsection{Interface utilisateur (dashboard)}

\subsubsection{Design et thème}

L'interface adopte le thème sombre \og GitHub Dark \fg{} (fonds \texttt{\#0d1117}, cartes \texttt{\#161b22}, bordures \texttt{\#30363d}, accent bleu \texttt{\#58a6ff}). Aucun framework CSS externe n'est utilisé : les styles sont en \texttt{<style>} inline dans le template \texttt{base.html}.

\subsubsection{Page dashboard}

La page d'accueil affiche :

\begin{itemize}
    \item \textbf{Statut des services} : tableau avec trois lignes (Simulateur, Pont MQTT-Kafka, Spark Inference). Chaque statut est coloré (vert pour \texttt{running}, rouge pour \texttt{stopped}). Ce fragment est rafraîchi automatiquement toutes les 5~secondes via HTMX.
    \item \textbf{Boutons de contrôle} : Démarrer, Arrêter, Redémarrer. Chaque bouton envoie une requête POST à la route correspondante via HTMX.
    \item \textbf{Dernière prédiction} : affiche le dernier score d'anomalie enregistré dans \texttt{alertes\_predictions}.
    \item \textbf{Machines} : liste de toutes les machines configurées avec leur statut (Actif/Inactif) et un bouton pour les activer ou désactiver directement depuis le dashboard.
    \item \textbf{Nombre de machines actives} : comptage depuis \texttt{machine\_config WHERE enabled = true}.
\end{itemize}

\subsubsection{Page configuration}

La page de configuration permet :

\begin{itemize}
    \item \textbf{Lister les machines} : tableau des machines configurées avec leurs paramètres (température cible, vibration cible, alpha) et leur état (Oui/Non). Chaque ligne dispose de boutons \og Activer/Désactiver \fg{} et \og Supprimer \fg{}.
    \item \textbf{Ajouter une machine} : formulaire avec ID, type (pipe/pump/generic), température et vibration cibles. L'insertion utilise \texttt{ON CONFLICT ... DO UPDATE}, ce qui permet aussi de modifier une machine existante.
    \item \textbf{Modifier le seuil d'alerte} : champ numérique avec valeur par défaut 0.7, limites 0.0--1.0, pas de 0.05.
\end{itemize}

\subsubsection{Accès à Grafana}

Un lien \og Grafana \fg{} dans la barre de navigation (onglet \texttt{target="\_blank"}) donne accès au dashboard Grafana à l'adresse \texttt{http://localhost:3000}.

\subsection{Vérification et tests}

Les vérifications suivantes ont été effectuées après déploiement :

\begin{itemize}
    \item \textbf{Endpoints HTTP} : les 3 pages (dashboard, config, status) répondent avec le code 200 et le HTML est valide.
    \item \textbf{Pipeline start} : les boutons Démarrer et Redémarrer lancent correctement le simulateur, le bridge, et soumettent le job Spark.
    \item \textbf{Modification du seuil} : la route \texttt{POST /config/threshold} met à jour la table \texttt{alert\_config} et retourne le nouveau seuil.
    \item \textbf{Gestion des machines} : l'activation/désactivation depuis l'interface web est prise en compte par le simulateur dans les 60~secondes. Les données des machines désactivées cessent d'être publiées dans Kafka.
    \item \textbf{Dockerisation} : les trois nouveaux conteneurs (simulateur, mqtt-bridge, webapp) démarrent sans erreur et les dépendances inter-conteneurs sont respectées.
    \item \textbf{Collecte des logs} : les conteneurs produisent des logs lisibles (timestamp, machine\_id, valeurs) facilitant le débogage.
\end{itemize}

\subsection{Améliorations possibles}

\begin{itemize}
    \item \textbf{Rechargement du seuil à chaud} : faire relire le seuil par Spark périodiquement (toutes les 60~secondes) via un thread séparé, pour éviter le redémarrage de la pipeline.
    \item \textbf{Authentification} : ajouter un login simple (FastAPI + JWT) pour sécuriser l'interface.
    \item \textbf{Historique des actions} : journaliser les actions (démarrage, arrêt, modification de config) dans une table \texttt{pipeline\_logs}.
    \item \textbf{Validation étendue} : vérifier que l'ID machine n'est pas vide et que les valeurs numériques sont dans des plages acceptables avant insertion.
    \item \textbf{Notifications} : envoyer un email ou webhook (Slack/Teams) lors du déclenchement d'une alerte.
\end{itemize}

\newpage

% ============================================================
% SYNTHESE POUR L'ENTRETIEN
% ============================================================
\section{Préparation à l'entretien : points clés à maîtriser}

\subsection{Questions potentielles et éléments de réponse}

\begin{enumerate}
    \item \textbf{Pourquoi Kafka entre MQTT et Spark ?}
    \textit{MQTT est un protocole publish/subscribe léger adapté aux capteurs, mais il ne garantit pas la persistance ni le rejeu. Kafka sert de tampon distribué : si Spark crashe, les messages s'accumulent dans Kafka et sont rejoués au redémarrage (\texttt{startingOffsets = "earliest"}).}

    \item \textbf{Pourquoi un seul modèle RandomForest pour température ET vibration ?}
    \textit{Plutôt que deux modèles séparés, on encode le type de capteur en feature catégorielle (StringIndexer). Cela permet au modèle d'apprendre des comportements spécifiques à chaque capteur tout en partageant la même structure. Avantage : maintenance simplifiée, un seul pipeline à déployer.}

    \item \textbf{Pourquoi des lags et des pentes plutôt que la valeur instantanée ?}
    \textit{Une valeur brute de 145\,°C est évidemment anormale, mais une montée de 95\,°C vers 100\,°C est déjà le signe d'une dérive. Les lags et la pente permettent une détection précoce : le modèle voit le changement avant que le seuil critique ne soit atteint. Gain : alerte 6 à 8~secondes plus tôt qu'un simple seuil fixe.}

    \item \textbf{Comment gérez-vous le déséquilibre des classes (90\% normal, 10\% anomalies) ?}
    \textit{Le RandomForest gère naturellement le déséquilibre via le vote majoritaire pondéré. Cependant, avec 30 arbres et depth=6 sur 10k échantillons, l'AUC de 0.80 pourrait être améliorée par du sur-échantillonnage ou du SMOTE.}

    \item \textbf{Quelle est la latence de bout en bout ?}
    \textit{Simulateur → MQTT → Kafka → Spark → prédiction → TimescaleDB : environ 2 à 4~secondes. Le trigger Spark est configuré à 2~secondes. Le goulot d'étranglement potentiel est le temps de traitement du micro-batch Spark (quelques centaines de millisecondes ici).}

    \item \textbf{Où est stocké le modèle et comment est-il chargé ?}
    \textit{Le modèle est sérialisé au format Spark ML (Parquet) dans \texttt{/home/spark/.ivy2/predictive\_model}. Il est chargé au démarrage du streaming via \texttt{PipelineModel.load()}. Le chemin doit être accessible depuis tous les exécuteurs Spark, d'où l'utilisation du volume partagé \texttt{spark\_ivy}.}
\end{enumerate}

\subsection{Schéma d'architecture pour l'entretien}

Pour l'entretien, pouvoir dessiner ce schéma sur un tableau ou papier :

\begin{center}
\small
\begin{tabular}{c}
\texttt{Simulateur (Python)} \\ $\downarrow$ \\ \texttt{MQTT (mosquitto)} \\ $\downarrow$ \\ \texttt{Pont MQTT→Kafka} \\ $\downarrow$ \\ \texttt{Kafka topic: sensor-data} \\ $\swarrow$ \quad $\downarrow$ \quad $\searrow$ \\
\texttt{TimescaleDB} \quad \texttt{MinIO (archive)} \quad \texttt{Spark ML (RandomForest)} \\
\qquad\qquad\qquad\qquad\qquad\qquad $\downarrow$ \\
\qquad\qquad\qquad\qquad\qquad\qquad \texttt{alerte\_predictions} \\
\qquad\qquad\qquad\qquad\qquad\qquad $\downarrow$ \\
\qquad\qquad\qquad\qquad\qquad\qquad \texttt{Grafana Dashboard} \\
\end{tabular}
\end{center}

\subsection{Points d'attention}

\begin{itemize}
    \item \textbf{Le modèle est unique} pour les deux capteurs (feature \texttt{sensor\_idx}).
    \item \textbf{Le dataset d'entraînement} a été généré par le même simulateur que le flux temps réel (cohérence de la distribution).
    \item \textbf{L'indicateur d'anomalie Grafana} est un \textit{stat panel} avec fond coloré, pas une jauge. La couleur (vert/jaune/rouge) donne le diagnostic ; la valeur brute du score (\textasciitilde0.026) n'est pas un pourcentage de risque mais la proportion d'arbres du RandomForest ayant voté \og anomalie \fg{} (voir explication section Visualisation Grafana).
    \item \textbf{Les checkpoints Spark} (\texttt{/tmp/checkpoint\_*}) doivent être nettoyés entre chaque redémarrage pour éviter des conflits d'état (\texttt{FileNotFoundException} sur les deltas).
    \item \textbf{NumPy} doit être accessible à l'exécuteur Spark (installé dans un répertoire du PYTHONPATH).
\end{itemize}

\newpage

\section*{Annexe : Structure des fichiers du projet}

\begin{lstlisting}[language=bash,caption={Arborescence du projet}]
raffinerie-iot/
  docker-compose.yml        # 11 services
  restart.sh                # docker compose up + verifications
  README.md
  LOGIQUE_SIMULATION.md     # Documentation physique

  # Scripts hotes (legacy, conserves pour historique)
  simulateur_capteurs.py    # Ancien simulateur (remplace par Docker)
  mqtt_to_kafka.py          # Ancien bridge (remplace par Docker)
  generate_training_data.py # Dataset labellise > MinIO
  train_model.py            # Entrainement du RandomForest

  spark/
    traitement_kpi.py       # Pipeline streaming + inference
    traitement_kpi_debug.py # Version debug
    train_model.py          # Entrainement
    predictive_model/       # Modele sauvegarde (3 stages)
    data/                   # Donnees d'entrainement

  simulateur/               # Conteneurise
    Dockerfile
    requirements.txt
    simulateur_capteurs.py  # Version multi-machines, lit la DB

  bridge/                   # Conteneurise
    Dockerfile
    requirements.txt
    mqtt_to_kafka.py        # Version wildcards, boucle d'attente Kafka

  webapp/                   # Interface web FastAPI + HTMX
    Dockerfile
    requirements.txt
    main.py
    templates/
      base.html
      dashboard.html
      config.html
      status_fragment.html

  initdb/                   # SQL d'initialisation TimescaleDB
    001_config_tables.sql   # Tables machine_config, alert_config

  volumes/
    mqtt/mosquitto.conf
\end{lstlisting}

\end{document}
